% =====================================================================
%  PHIEU KHAO SAT HOC SINH  --  tep DOC LAP, bien dich rieng de in phat
%  Bien dich: xelatex phieu_khao_sat.tex
%
%  Dung 2 lan:  - TRUOC khi ap dung (dau nam)  -> so lieu cho Bang 2
%               - SAU khi ap dung  (cuoi dot)  -> so lieu cho Bang 10
%  Cung mot phieu, chi doi o "Thoi diem khao sat".
% =====================================================================
\documentclass[12pt,a4paper]{article}
\usepackage{fontspec}
\usepackage[left=2.2cm,right=1.8cm,top=2cm,bottom=2cm]{geometry}
\usepackage{array,longtable,enumitem,multirow}
\usepackage{amssymb}   % \square cho o danh dau
\setmainfont{Times New Roman}
% Neu may bao khong tim thay font, doi thanh:  \setmainfont{TeX Gyre Termes}
\setlength{\parindent}{0pt}
\renewcommand{\arraystretch}{1.5}
\pagestyle{empty}

\begin{document}

\begin{center}
{\large\textbf{PHIẾU KHẢO SÁT Ý KIẾN HỌC SINH}}\\[2pt]
{\footnotesize Về việc sử dụng hệ thống Ngân hàng đề trực tuyến trong học tập môn Toán}
\end{center}

\vspace{4pt}
\textit{\small Phiếu này \textbf{không ghi tên}. Kết quả chỉ dùng để cải tiến
việc dạy học, không ảnh hưởng đến điểm số của em. Em trả lời đúng suy nghĩ thật
của mình --- đó mới là điều có ích.}

\vspace{6pt}
Lớp: \dotfill\ \ Thời điểm khảo sát: $\square$ Đầu năm \quad $\square$ Cuối đợt

\vspace{8pt}
\textbf{Phần A. Em đánh dấu $\times$ vào ô phù hợp nhất.}

\vspace{2pt}
{\footnotesize 1 --- Hoàn toàn không đồng ý \quad 2 --- Không đồng ý \quad
3 --- Bình thường \quad 4 --- Đồng ý \quad 5 --- Hoàn toàn đồng ý}

\vspace{4pt}
\begin{longtable}{|p{0.7cm}|p{10.6cm}|c|c|c|c|c|}
\hline
\textbf{TT} & \textbf{Nội dung} & \textbf{1} & \textbf{2} & \textbf{3} & \textbf{4} & \textbf{5} \\ \hline
\endhead

\multicolumn{7}{|l|}{\textbf{I. Thái độ với môn Toán}} \\ \hline
1 & Em thấy hứng thú khi học môn Toán & & & & & \\ \hline
2 & Em mong đến tiết Toán & & & & & \\ \hline
3 & Em thấy việc luyện tập Toán là nhẹ nhàng, không áp lực & & & & & \\ \hline

\multicolumn{7}{|l|}{\textbf{II. Mức độ tự tin}} \\ \hline
4 & Em tự tin khi làm bài kiểm tra Toán & & & & & \\ \hline
5 & Gặp bài chưa biết làm, em tin mình sẽ tìm ra cách & & & & & \\ \hline
6 & Em nắm được mình đang yếu ở phần kiến thức nào & & & & & \\ \hline

\multicolumn{7}{|l|}{\textbf{III. Thói quen tự luyện tập}} \\ \hline
7 & Em có luyện thêm bài tập Toán ngoài bài thầy/cô giao & & & & & \\ \hline
8 & Em tìm được đủ bài tập đúng dạng mình cần luyện & & & & & \\ \hline
9 & Khi làm sai, em có quay lại làm lại dạng bài đó & & & & & \\ \hline

\multicolumn{7}{|l|}{\textbf{IV. Về hệ thống Ngân hàng đề} \textit{(chỉ trả lời ở lần khảo sát cuối đợt)}} \\ \hline
10 & Hệ thống dễ dùng, em không phải hỏi ai cách sử dụng & & & & & \\ \hline
11 & Em thích việc mỗi lần tạo đề lại được một đề khác số liệu & & & & & \\ \hline
12 & Được chấm điểm ngay giúp em biết mình sai ở đâu & & & & & \\ \hline
13 & Phần thầy/cô AI giảng lại giúp em hiểu bài hơn & & & & & \\ \hline
14 & Em tin vào lời giải mà hệ thống đưa ra & & & & & \\ \hline
15 & Em muốn tiếp tục dùng hệ thống ở học kì sau & & & & & \\ \hline
16 & Em sẽ giới thiệu hệ thống cho bạn ở lớp khác & & & & & \\ \hline
\end{longtable}

\vspace{4pt}
\textbf{Phần B. Em trả lời ngắn gọn.}

\vspace{4pt}
17. Trung bình mỗi tuần em vào hệ thống khoảng bao nhiêu lần?

\quad $\square$ Không vào \quad $\square$ 1--2 lần \quad $\square$ 3--5 lần
\quad $\square$ Trên 5 lần

\vspace{6pt}
18. Chức năng em dùng nhiều nhất \textit{(chọn tối đa hai)}:

\quad $\square$ Tạo đề luyện tập \quad $\square$ Làm bài, chấm điểm
\quad $\square$ Xem lời giải \quad $\square$ Hỏi thầy/cô AI

\vspace{6pt}
19. Điều em thấy có ích nhất:

\vspace{2pt}\dotfill\\[4pt]\dotfill

\vspace{6pt}
20. Điều em thấy còn khó chịu, hoặc mong được cải tiến:

\vspace{2pt}\dotfill\\[4pt]\dotfill

\vspace{10pt}
\begin{center}\textit{Cảm ơn em đã trả lời.}\end{center}

\clearpage

% =====================================================================
%  TRANG NAY DANH CHO GIAO VIEN -- KHONG IN PHAT CHO HOC SINH
% =====================================================================
\begin{center}{\large\textbf{HƯỚNG DẪN XỬ LÍ PHIẾU}}\\{\footnotesize (trang dành cho giáo viên --- không in phát)}\end{center}

\vspace{6pt}
\textbf{1. Nhập liệu.} Mỗi phiếu là một dòng. Cột: \texttt{ID}, \texttt{Lop},
\texttt{ThoiDiem} (1 = đầu năm, 2 = cuối đợt), rồi \texttt{C1} đến \texttt{C16}
nhận giá trị 1--5. Câu 10--16 bỏ trống ở lần khảo sát đầu năm.

\textbf{2. Kiểm tra độ tin cậy trước khi dùng.} Tính hệ số Cronbach's $\alpha$
cho \textbf{từng nhóm} I, II, III, IV --- không tính chung cả 16 câu. Đạt từ
0{,}7 trở lên là dùng được; dưới 0{,}6 thì xem lại có câu nào học sinh hiểu
nhầm và loại câu đó ra.

\textbf{3. So sánh đầu năm với cuối đợt.} Dùng kiểm định t cho hai mẫu ghép cặp
trên từng nhóm I, II, III. Nhóm IV chỉ có một lần đo nên chỉ báo điểm trung
bình.

\textbf{4. Đối chiếu với nhật kí hệ thống.} Đây là bước làm bài viết mạnh lên
hẳn. Câu 17 là học sinh \textit{tự khai}; số lần vào thật thì hệ thống có ghi.
Nếu hai nguồn cùng chỉ một hướng, kết luận vững hơn nhiều so với chỉ dựa vào
phiếu hỏi.

\vspace{6pt}
\textbf{Ba điều cần nói rõ trong bài viết:}

\begin{enumerate}[leftmargin=1.1cm, itemsep=3pt]
  \item Phiếu \textbf{ẩn danh hoàn toàn} --- đã ghi ngay đầu phiếu.
  \item Phát phiếu khi \textbf{không có mặt giáo viên} hoặc để học sinh tự bỏ
        vào hộp kín, nhằm giảm việc các em trả lời cho vừa lòng thầy cô.
  \item \textbf{Thừa nhận hạn chế:} khảo sát tự báo cáo luôn có thiên lệch, nên
        phiếu chỉ là một trong ba nguồn --- cùng với điểm kiểm tra và nhật kí
        hệ thống. Hội đồng đánh giá cao việc tác giả tự nêu hạn chế hơn là giấu
        đi.
\end{enumerate}

\end{document}
